\chapter{Análisis de costes}\label{cap:Costes}
% Índice v4, cap. 5: coste real con IA frente a estimación clásica con el mismo
% alcance. Datos reales del apartado 4.5; la estimación clásica explica su método.

\section{Costes del proyecto con IA}\label{sec:CostesIA}
% Personal (horas por fase, tarifa de referencia citada); herramientas de IA
% (suscripción de Claude para el desarrollo y API con gasto propio mínimo para el
% RAG); infraestructura (local); equipo y licencias.
\pendiente{costes reales del desarrollo}

\section{Infraestructura LLM: opciones y costes}\label{sec:InfraLLM}
% Tutoría anotada el 2026-11-09. API de Claude con gasto propio mínimo y límite de
% unos 10 EUR en la consola; generación con Haiku; embeddings con Voyage AI (cuota
% gratuita) o con un modelo local; modelo local de generación como alternativa.
% Coste real de la demo y de la evaluación del RAG; precios con fuente y fecha.
\pendiente{opciones y costes de la infraestructura LLM}

\section{Coste de operación}\label{sec:CosteOperacion}
% Coste mensual en producción (alojamiento, LLM y embeddings por consulta,
% almacenamiento, monitorización), también por agricultor y mes.
\pendiente{coste de operación}

\section{Estimación con el enfoque clásico}\label{sec:EstimacionClasica}
% Equipo tipo sin IA con tarifas citadas; desglose por fase contrastado con
% COCOMO II o puntos de función; mismo alcance.
\pendiente{estimación clásica}

\section{Comparativa: enfoque clásico frente a IA}\label{sec:Comparativa}
% Tabla por fase y barras agrupadas; costes ocultos de la IA; calidad.
\pendiente{comparativa}

\section{Análisis de sensibilidad}\label{sec:Sensibilidad}
% Tarifa, precio de tokens y tamaño; punto de equilibrio; limitaciones.
\pendiente{sensibilidad}
